
\documentclass[aspectratio=169]{beamer}
\usetheme{Madrid}
\usecolortheme{default}
\usepackage[utf8]{inputenc}
\usepackage{amsmath, amssymb}
\usepackage{booktabs}
\usepackage{ragged2e}
\usepackage{multicol}

\title{GED Math Unit: Decimals and Fractions}
\author{GED Preparation}
\date{\today}

\begin{document}

\begin{frame}
  \titlepage
\end{frame}

\begin{frame}{Unit Overview}
\tableofcontents
\end{frame}

% =============================
\section{Lesson 1: Understanding and Comparing Decimals}
% =============================
\begin{frame}{Lesson 1 Objectives}
\begin{itemize}
  \item Understand the place value of decimals (tenths, hundredths, thousandths).
  \item Compare and order decimals.
  \item Convert between decimals and fractions.
\end{itemize}
\end{frame}

\begin{frame}{Lesson 1 --- Key Ideas}
\begin{itemize}
  \item Place value chart: $1.2\underline{3}\,\underline{4}$ has 3 in the hundredths place and 4 in the thousandths place.
  \item Compare decimals using place value and by aligning digits with zeros as placeholders.
  \item Convert decimals to fractions by reading the place value (\textit{e.g.,} $0.45=\frac{45}{100}$).
\end{itemize}
\end{frame}

\begin{frame}{Lesson 1 --- Worked Examples}
\begin{enumerate}
  \item Write $0.45$ in words: \textbf{forty-five hundredths}.
  \item Which is greater: $0.7$ or $0.65$? \quad Compare tenths: $7$ tenths vs $6$ tenths $\Rightarrow$ \textbf{$0.7>0.65$}.
  \item Convert $0.25$ to a fraction: $0.25=\frac{25}{100}=\frac{1}{4}$.
  \item Convert $\frac{3}{10}$ to a decimal: \textbf{$0.3$}.
  \item Order $0.56$, $0.5$, $0.506$ by rewriting as $0.560$, $0.500$, $0.506$ $\Rightarrow$ \textbf{$0.5$, $0.506$, $0.56$}.
\end{enumerate}
\end{frame}

% =============================
\section{Lesson 2: Operations with Decimals}
% =============================
\begin{frame}{Lesson 2 Objectives}
\begin{itemize}
  \item Add and subtract decimals by aligning decimal points.
  \item Multiply and divide decimals by tracking decimal places.
\end{itemize}
\end{frame}

\begin{frame}{Lesson 2 --- Strategy Notes}
\begin{itemize}
  \item \textbf{Addition/Subtraction:} line up the decimal points; use zeros as placeholders.
  \item \textbf{Multiplication:} multiply as whole numbers, then place the decimal; the total decimal places equals the sum from the factors.
  \item \textbf{Division:} move the decimal in the divisor to make it a whole number and move the dividend's decimal the same amount.
\end{itemize}
\end{frame}

\begin{frame}{Lesson 2 --- Worked Examples}
\begin{enumerate}
  \item $4.52+3.6=4.52+3.60=\textbf{8.12}$.
  \item $5.7-2.45=5.70-2.45=\textbf{3.25}$.
  \item $0.4\times0.6=4\times6=24$ with two decimal places $\Rightarrow$ \textbf{$0.24$}.
  \item $6.3\div0.9 \Rightarrow 63\div9=\textbf{7}$.
  \item $0.25\times12=25\times12=300$ with two decimal places $\Rightarrow$ \textbf{$3.00$}.
\end{enumerate}
\end{frame}

% =============================
\section{Lesson 3: Understanding Fractions}
% =============================
\begin{frame}{Lesson 3 Objectives}
\begin{itemize}
  \item Identify numerator and denominator.
  \item Simplify fractions using greatest common factor (GCF).
  \item Find equivalent fractions.
\end{itemize}
\end{frame}

\begin{frame}{Lesson 3 --- Key Ideas}
\begin{itemize}
  \item A fraction $\frac{a}{b}$ represents $a$ parts out of $b$ equal parts ($b\neq 0$).
  \item To simplify, divide numerator and denominator by their GCF.
  \item Create an equivalent fraction by multiplying (or dividing) top and bottom by the same nonzero number.
\end{itemize}
\end{frame}

\begin{frame}{Lesson 3 --- Worked Examples}
\begin{enumerate}
  \item Parts of $\frac{3}{5}$: numerator $=3$, denominator $=5$.
  \item Simplify $\frac{8}{12}$: GCF$=4 \Rightarrow \frac{8}{12}=\frac{2}{3}$.
  \item Equivalent to $\frac{1}{2}$ with denominator $10$: $\frac{1}{2}=\frac{5}{10}$.
  \item Simplify $\frac{45}{60}$: GCF$=15 \Rightarrow \frac{45}{60}=\frac{3}{4}$.
  \item Are $\frac{4}{8}$ and $\frac{1}{2}$ equivalent? Yes, $\frac{4}{8}\to\frac{1}{2}$.
\end{enumerate}
\end{frame}

% =============================
\section{Lesson 4: Operations with Fractions}
% =============================
\begin{frame}{Lesson 4 Objectives}
\begin{itemize}
  \item Add and subtract with common denominators (or find the LCD).
  \item Multiply numerators and denominators; simplify.
  \item Divide by multiplying by the reciprocal.
\end{itemize}
\end{frame}

\begin{frame}{Lesson 4 --- Strategy Notes}
\begin{itemize}
  \item \textbf{Addition/Subtraction:} Use the least common denominator (LCD); convert, then add/subtract numerators.
  \item \textbf{Multiplication:} Multiply across; reduce before or after.
  \item \textbf{Division:} Keep-Change-Flip: keep the first fraction, change to multiply, flip the second.
\end{itemize}
\end{frame}

\begin{frame}{Lesson 4 --- Worked Examples}
\begin{enumerate}
  \item $\frac{1}{4}+\frac{2}{4}=\frac{3}{4}$.
  \item $\frac{3}{8}-\frac{1}{8}=\frac{2}{8}=\frac{1}{4}$.
  \item $\frac{2}{3}\times\frac{3}{4}=\frac{6}{12}=\frac{1}{2}$.
  \item $\frac{5}{6}\div\frac{2}{3}=\frac{5}{6}\times\frac{3}{2}=\frac{15}{12}=\mathbf{1\frac{1}{4}}$.
  \item $\frac{1}{2}+\frac{1}{3}$: LCD$=6 \Rightarrow \frac{3}{6}+\frac{2}{6}=\frac{5}{6}$.
\end{enumerate}
\end{frame}

% =============================
\section{Lesson 5: Converting Between Fractions and Decimals}
% =============================
\begin{frame}{Lesson 5 Objectives}
\begin{itemize}
  \item Convert fractions to decimals via long division.
  \item Convert terminating and repeating decimals to fractions.
\end{itemize}
\end{frame}

\begin{frame}{Lesson 5 --- Key Ideas}
\begin{itemize}
  \item Fractions with denominators made of $2$s and $5$s (like $8$, $20$, $125$) produce terminating decimals.
  \item Repeating decimals represent rational numbers: $0.\overline{3}=\frac{1}{3}$.
  \item To write a decimal as a fraction, read its place value and simplify.
\end{itemize}
\end{frame}

\begin{frame}{Lesson 5 --- Worked Examples}
\begin{enumerate}
  \item $\frac{1}{4}\Rightarrow 1\div4=0.25$.
  \item $\frac{2}{5}\Rightarrow 2\div5=0.4$.
  \item $0.75=\frac{75}{100}=\frac{3}{4}$.
  \item $0.\overline{3}=\frac{1}{3}$ (standard result for repeating $3$).
  \item $0.875=\frac{875}{1000}=\frac{7}{8}$.
\end{enumerate}
\end{frame}

% =============================
\section{Lesson 6: Applications of Decimals and Fractions}
% =============================
\begin{frame}{Lesson 6 Objectives}
\begin{itemize}
  \item Solve real-world word problems using decimals and fractions.
  \item Interpret and compute with measurements, money, and ratios.
\end{itemize}
\end{frame}

\begin{frame}{Lesson 6 --- Worked Examples}
\begin{enumerate}
  \item A shirt costs \$19.99 and tax is \$1.60. Total cost: \$19.99+\$1.60=\textbf{\$21.59}.
  \item You ran $\frac{3}{4}$ mile and your friend ran $\frac{1}{2}$ mile. Difference: $\frac{3}{4}-\frac{1}{2}=\frac{1}{4}$ mile.
  \item Find $0.6$ of $50$: $0.6\times50=\textbf{30}$.
  \item Recipe uses $\frac{1}{2}$ cup sugar and $\frac{1}{4}$ cup flour. Total: $\frac{1}{2}+\frac{1}{4}=\frac{3}{4}$ cup.
  \item You have \$5 and spend \$3.25. Remaining: \$5.00-\$3.25=\textbf{\$1.75}.
\end{enumerate}
\end{frame}

% =============================
\section{Exit Ticket \& Practice}
% =============================
\begin{frame}{Quick Check (Exit Ticket)}
\begin{enumerate}
  \item Write $0.306$ as a fraction in simplest form.
  \item Order: $0.68,\;0.6,\;0.605$.
  \item Compute: $3.75-1.8$.
  \item Simplify: $\frac{18}{24}$.
  \item Compute: $\frac{2}{5}+\frac{3}{10}$.
\end{enumerate}
\end{frame}

\begin{frame}{Guided Practice (Optional)}
\begin{multicols}{2}
\begin{enumerate}
  \item $0.07+0.6$
  \item $1.2\times0.3$
  \item $\frac{4}{9}\times\frac{3}{8}$
  \item $\frac{7}{10}\div\frac{1}{5}$
  \item Convert $0.2$ to a fraction
  \item Convert $\frac{9}{20}$ to decimal
\end{enumerate}
\end{multicols}
\end{frame}

\begin{frame}{Summary}
\begin{itemize}
  \item Decimals and fractions represent parts of a whole and convert between each other.
  \item Operations follow consistent rules: align decimals; use LCDs; multiply across; divide by reciprocal.
  \item Apply these skills to money, measurement, and data contexts on the GED.
\end{itemize}
\end{frame}

\end{document}
